\section*{Введение}
\addcontentsline{toc}{section}{Введение}

Система доменных имён (DNS) является одним из фундаментальных компонентов современных компьютерных сетей.
Она обеспечивает преобразование символьных доменных имён в числовые IP-адреса, необходимые для маршрутизации пакетов.
Без DNS пользователь был бы вынужден запоминать IP-адреса каждого сетевого ресурса вместо привычных имён типа \texttt{example.com}.

В рамках данной практической работы рассматривается развёртывание собственного DNS-сервера на базе пакета \textbf{BIND9}
(Berkeley Internet Name Domain) в операционной системе Ubuntu~20.04~LTS~\cite{bind9_admin}.

Целью работы является настройка полноценного DNS-сервера для локальной сети, включающая:
\begin{itemize}
  \item создание прямой зоны (forward zone) для преобразования имён в IP-адреса;
  \item создание обратной зоны (reverse zone) для обратного преобразования;
  \item настройку DDNS для автоматической регистрации клиентов через DHCP.
\end{itemize}

Стенд: виртуальная машина \texttt{gateway} (VirtualBox, Ubuntu~20.04),
интерфейсы \texttt{enp0s3}~(WAN/NAT) и \texttt{enp0s8}~(LAN,~\texttt{192.168.\cfgLabStudentN.1/24}).
